处理 PDE 方程组 (而非单个方程) 时, 一般不着重采用限制更强的椭圆性定义. 也就是说, 仍将椭圆方程组定义为没有实特征的方程组, 而特征由使用适当权重定义的算子主部确定. 不过, 读者应了解变分法中出现的一些特别重要的椭圆方程组例子. 设 $\Omega$ 为 $\mathbb{R}^n$ 中的区域, $\bm{u}:\Omega\to\mathbb{R}^N$. 考虑由 $N$ 个二阶散度形式微分方程组成的方程组
\begin{equation}\label{eq:second-order-divergence-system-definition}
\frac{\partial}{\partial x_k}\left(A_{kl}^{IJ}(\bm{x})\frac{\partial u_J}{\partial x_l}\right)+c^{IJ}(\bm{x})u_J(\bm{x})=0,\quad I=1,\ldots,N.
\end{equation}


这里假设系数 $A_{kl}^{IJ}$ 和 $c^{IJ}$, $I,J=1,\ldots,N$, $k,l=1,\ldots,n$, 足够光滑, 并约定: 重复的大写指标从 $1$ 到 $N$ 求和, 重复的小写指标从 $1$ 到 $n$ 求和.

注意, 若对每个 $\bm{x}\in\Omega$ 以及每个非零 $N\times n$ 矩阵 $\bm{M}$ 都有
\begin{equation}\label{eq:elliptic-system-full-matrix-positivity}
M_k^IA_{kl}^{IJ}M_l^J>0
\end{equation}
则只需取 ``显然的\textquotesingle\textquotesingle\ 主部 (赋予每个方程及因变量权重 $1$), 就能证明该方程组椭圆. 条件 \eqref{eq:elliptic-system-full-matrix-positivity}, 或其一致形式
\begin{equation}\label{eq:elliptic-system-full-matrix-uniform-bound}
M_k^IA_{kl}^{IJ}M_l^J\geq\theta|\bm{M}|^2
\end{equation}
对每个 $\bm{x}\in\Omega$ 和每个非零 $N\times n$ 矩阵 $\bm{M}$ 以某个 $\theta>0$ 成立, 常被用作椭圆方程组的定义. 然而, 这样的定义并不适用于 Stokes 方程组等方程组.

\PDEIndexMark{idx:legendre-hadamard-condition-and-terminology-caveat}另一个重要的 ``椭圆性条件\textquotesingle\textquotesingle\ 是 \textbf{Legendre-Hadamard} 条件
\begin{equation}\label{eq:legendre-hadamard-positivity-condition}
\eta^I\xi_kA_{kl}^{IJ}\eta^J\xi_l>0
\end{equation}
它要求对每个 $\bm{x}\in\Omega$ 及每个非零 $\bm{\eta}\in\mathbb{R}^N$ 和 $\bm{\xi}\in\mathbb{R}^n$ 都成立. 一致形式是指存在 $\theta>0$, 使得对每个 $\bm{x}\in\Omega$ 及每个非零 $\bm{\eta}\in\mathbb{R}^N$ 和 $\bm{\xi}\in\mathbb{R}^n$ 均有
\begin{equation}\label{eq:legendre-hadamard-uniform-condition}
\eta^I\xi_kA_{kl}^{IJ}\eta^J\xi_l>\theta|\bm{\eta}|^2|\bm{\xi}|^2
\end{equation}
对于弹性理论中的许多问题, 这些条件在物理上比 \eqref{eq:elliptic-system-full-matrix-positivity} 或 \eqref{eq:elliptic-system-full-matrix-uniform-bound} 更合理. 注意, \eqref{eq:legendre-hadamard-positivity-condition} 和 \eqref{eq:legendre-hadamard-uniform-condition} 远弱于相应的条件 \eqref{eq:elliptic-system-full-matrix-positivity} 和 \eqref{eq:elliptic-system-full-matrix-uniform-bound}. (这些不等式只须对秩为 $1$ 的 $N\times n$ 矩阵成立.) 尽管如此, \eqref{eq:legendre-hadamard-positivity-condition} 和 \eqref{eq:legendre-hadamard-uniform-condition} 有时仍称为\emph{强椭圆性}条件. 这个例子提醒读者, 围绕椭圆方程组的术语未必合乎情理. 更重要的是, 这些定义并无普遍一致的约定. 阅读文献时, 必须仔细留意不同作者使用的定义.
